\setchoices{label=(\Alph*)}

\begin{question}[index=1]
设函数 $f(x)=\displaystyle\int_0^x\eu^{t^2}\dd t$，对任意的 $x\ne0$，存在介于 $0$ 和 $x$ 之间的 $\xi$，有 $f(\xi)=\dfrac{f(x)+f(0)}2$，则 $\displaystyle\lim_{x\to0}\frac{\xi}{x}=$
\begin{choices}[columns=4]
\item $0$
\item $\dfrac13$
\item $\dfrac12$
\item $1$
\end{choices}
\end{question}
\QuestionBreak

\begin{question}[index=2]
设函数 $f(x,y)=\begin{cases}1,&0<y<x^2,\\0,&\text{其他},\end{cases}$ 则下列关于 $f(x,y)$ 在点 $(0,0)$ 的说法正确的有几个？
\par\circlednumber{1}\quad 连续
\par\circlednumber{2}\quad 偏导数存在
\par\circlednumber{3}\quad 可微
\par\circlednumber{4}\quad 沿任意方向的方向导数存在
\begin{choices}[columns=4]
\item $1$
\item $2$
\item $3$
\item $4$
\end{choices}
\end{question}
\QuestionBreak

\begin{question}[index=3]
区域 $D$ 由 $r=1-\cos\theta$（$0\le\theta\le\pi$）以及 $\theta=\pi$ 所围，则 $\displaystyle\iint_D(x+y)\dd\sigma=$
\begin{choices}[columns=4]
\item $\dfrac23+\dfrac{\pi}{8}$
\item $\dfrac43-\dfrac58\pi$
\item $0$
\item $\dfrac73-\dfrac38\pi$
\end{choices}
\end{question}
\QuestionBreak

\begin{question}[index=4]
设在 $[0,1]$ 上连续的函数 $f(x)\ge0$。若曲线 $y=f(x)$ 在 $[0,1]$ 上构成的\allowbreak\mbox{曲边梯形}面积为 $1$，形心为 $\left(\dfrac12,\dfrac12\right)$，则该\allowbreak\mbox{曲边梯形}分别绕 $x$ 轴和 $y$ 轴旋转一周所得旋转体的体积之和为
\begin{choices}[columns=4]
\item $\pi$
\item $2\pi$
\item $1$
\item $2$
\end{choices}
\end{question}
\QuestionBreak

\begin{question}[index=5]
设 $3$ 阶方阵 $A,B$ 满足方程 $A-2B^2=AB$，且 $B$ 可逆，则下列等式不正确的是
\begin{choices}[columns=2]
\item $AB^{-1}=2B+A$
\item $(B^{-1}-E)^{-1}=\dfrac12B^{-1}A$
\item $(E-B)^{-1}=\dfrac12(A+2E)$
\item $BA^{-1}B=\dfrac12(E-B)$
\end{choices}
\end{question}
\QuestionBreak

\begin{question}[index=6]
设 $\alpha_1,\alpha_2,\beta_1,\beta_2$ 为 $3$ 维向量组，且 $\alpha_1,\alpha_2$ 线性无关，$\beta_1,\beta_2$ 也线性无关，则下列说法正确的是
\begin{choices}[columns=1]
\item $\alpha_1,\alpha_2$ 与 $\beta_1,\beta_2$ 是等价向量组
\item 方程组 $\mat{\alpha_1^{\mathrm T}\\\alpha_2^{\mathrm T}}x=0$ 与方程组 $\mat{\beta_1^{\mathrm T}\\\beta_2^{\mathrm T}}x=0$ 同解
\item 方程组 $\mat{\alpha_1^{\mathrm T}\\\alpha_2^{\mathrm T}}x=0$ 与方程组 $\mat{\beta_1^{\mathrm T}\\\beta_2^{\mathrm T}}x=0$ 必有非零公共解
\item 存在非零向量 $\xi$ 可由 $\alpha_1,\alpha_2$ 线性表示，也可由 $\beta_1,\beta_2$ 线性表示
\end{choices}
\end{question}
\QuestionBreak

\begin{question}[index=7]
下列矩阵与 $A=\mat{1&0&1\\0&2&2\\0&0&2}$ 相似的是
\begin{choices}[columns=4]
\item $\mat{2&1&2\\0&1&-2\\0&0&2}$
\item $\mat{1&2&0\\1&3&0\\0&0&2}$
\item $\mat{1&0&1\\1&2&1\\0&0&2}$
\item $\mat{2&1&1\\1&3&0\\2&0&0}$
\end{choices}
\end{question}
\QuestionBreak

\begin{question}[index=8]
设随机变量 $X$ 的概率分布为 $X\dist\left(\begin{array}{ccc}0&1&2\\a&b&c\end{array}\right)$，其中 $a,b,c$ 均为常数。期望 $E(X)=\dfrac23$，又 $n$ 维向量 $\alpha_1,\alpha_2,\alpha_3$ 线性无关。若 $\alpha_1+\alpha_2+\alpha_3$，$\alpha_2+X\alpha_3$，$X\alpha_1-2\alpha_3$ 线性相关的概率为 $\dfrac16$，则 $(a,b,c)=$
\begin{choices}[columns=4]
\item $\left(\dfrac12,\dfrac16,\dfrac13\right)$
\item $\left(\dfrac12,\dfrac13,\dfrac16\right)$
\item $\left(\dfrac16,\dfrac16,\dfrac23\right)$
\item $\left(\dfrac{13}{24},\dfrac14,\dfrac5{24}\right)$
\end{choices}
\end{question}
\QuestionBreak

\begin{question}[index=9]
设随机变量 $X_1,X_2,\cdots,X_n,\cdots$ 相互独立同分布，如果分布函数为 $F(x)=\dfrac12+\dfrac1\pi\arctan\dfrac{x}{B}$，$-\infty<x<+\infty$，其中 $B$ 为大于 $0$ 的常数，则此序列\quad（\quad）辛钦大数定律条件。
\begin{choices}[columns=2]
\item 满足
\item $B$ 取适当值时满足
\item 不满足
\item 无法判断
\end{choices}
\end{question}
\QuestionBreak

\begin{question}[index=10]
设总体 $X\dist N(\mu,1)$，$X_1,\cdots,X_n$ 为来自总体 $X$ 的简单随机样本。假设检验 \allowbreak\mbox{$H_0\colon \mu\le0$}，\allowbreak\mbox{$H_1\colon\mu>0$}，取拒绝域 $W=\{\overline X\ge c\}$，其中 $\overline X$ 为样本均值。那么对固定的样本容量 $n$，当 $\mu=1$ 时，\allowbreak\mbox{犯第二类错误}的概率 $\beta$
\begin{choices}[columns=2]
\item 随 $c$ 的增大而减少
\item 随 $c$ 的增大而增大
\item 随 $c$ 的增大而保持不变
\item 无法判定
\end{choices}
\end{question}
